% NOTE: numbers in this file are placeholders until the formal batch is analysed.
% The reporting rule of this project: every number in the paper must be traceable
% to a raw batch under results/, and no unfavourable function may be dropped.
Traditional swarm intelligence algorithms evaluate their entire population at
every generation. When one objective evaluation consumes a costly simulation,
this full-population paradigm spends most of the budget on candidates whose
fitness is already implied by their neighbours. This paper studies the resulting
resource-allocation problem for the grey wolf optimizer (GWO), a leader-driven
algorithm whose search direction is controlled by only three positions. We
propose the Resource-Aware Sparse-Evaluation Grey Wolf Optimizer (RASE-GWO),
which decouples candidate generation from candidate evaluation. A virtual
population of $M$ wolves moves under the unmodified GWO equations at negligible
cost, but only $K \ll M$ of them consume true evaluations; the quota is assigned
by a dual-potential geometric screening rule that selects candidates near the
leader-consensus centre for refinement and candidates near the minority leader
to keep a second basin alive. Only verified positions may enter the
$\alpha/\beta/\delta$ hierarchy, which prevents unevaluated movement from
contaminating the search direction, and a tabular Q-learning agent reallocates
the split between the two candidate groups from a bounded, scale-invariant
reward. On the CEC2017 benchmark at $D=50$ with an identical budget of
$10^{5}$ true evaluations and an identical number of generations, RASE-GWO
improves the mean error of canonical GWO on LIBRARY-NUMBERS-PLACEHOLDER of 29
functions with a median error ratio of RATIO-PLACEHOLDER, while covering the
same number of true evaluations. Ablation studies show that the geometric
screening, the over-sampled candidate pool and the verification decoupling each
contribute, and that evaluating wolves that are far from the swarm centroid --
the criterion inherited from sparse-evaluation particle swarm frameworks --
converts very few evaluations into improvements. We also report two negative
results that delimit the framework: for GWO the sparse-evaluation schedule does
not pay off when the per-round quota is reduced without over-sampling, and a
naive population reduction is clearly inferior. The additional cost of
RASE-GWO is arithmetic rather than evaluative, which makes it attractive for
optimization driven by expensive simulations.
